%% ma: L12-B15-0015
%% lop: 12
%% chuong: 5
%% bai: 15
%% dang: 12.15.12
%% loai: TN4
%% muc: VD
%% dap_an: "B"
%% nguon: "(NBV)-10 ĐỀ KIỂM TRA HK2-TOÁN 12-DUNG CHUNG 3 SÁCH.pdf (H01, Đề số 1), Phần 1 Câu 6"
%% kien_thuc: "Bài 15 (đường thẳng trong không gian, góc giữa hai đường thẳng), vectơ đơn vị và đường phân giác"
%% trang_thai: chinh-thuc
%% ghi_chu: "Đáp án gốc B đúng (giải độc lập: hai đường thẳng cắt nhau tại A(1;-3;5); u.v = -5 < 0 nên phân giác góc nhọn có vectơ chỉ phương u/|u| - v/|v| ∥ (2;-5;11)). Đã duyệt 19/10/2026."
\begin{ex}
Trong không gian $Oxyz$, cho đường thẳng $d\colon\begin{cases}x=1+3t\\y=-3\\z=5+4t\end{cases}$; $\Delta$ là đường thẳng qua điểm $A(1;-3;5)$ và có vectơ chỉ phương $\vec u=(1;2;-2)$. Đường phân giác của góc nhọn tạo bởi $d$ và $\Delta$ có phương trình là
\choice{$\begin{cases}x=-1+2t\\y=2-5t\\z=6+11t\end{cases}$}{\True $\begin{cases}x=1+2t\\y=-3-5t\\z=5+11t\end{cases}$}{$\begin{cases}x=1+7t\\y=-3+5t\\z=5+t\end{cases}$}{$\begin{cases}x=1-t\\y=-3\\z=5+7t\end{cases}$}
\loigiai{Điểm $A$ thuộc cả $d$ (ứng với $t=0$) và $\Delta$ nên hai đường thẳng cắt nhau tại $A(1;-3;5)$. Vectơ chỉ phương của $d$ là $\vec v=(3;0;4)$.

Vì $\vec u\cdot\vec v=3+0-8=-5<0$ nên góc giữa $\vec u$ và $\vec v$ là góc tù. Do đó phân giác của góc nhọn có vectơ chỉ phương $\dfrac{\vec u}{|\vec u|}-\dfrac{\vec v}{|\vec v|}=\dfrac13(1;2;-2)-\dfrac15(3;0;4)=\dfrac2{15}(-2;5;-11)$, tức $(2;-5;11)$.

Phương trình đường phân giác là $\begin{cases}x=1+2t\\y=-3-5t\\z=5+11t.\end{cases}$ Chọn B.}
\end{ex}
